\section{Lessons learned}
\label{sec:lessons}

\subsection{What worked}

\paragraph{Pydantic-driven JSON Schemas as the contract between backend and frontend.}
Every strategy in \texttt{backend/strategies/} declares a Pydantic \texttt{Config}
class; \texttt{GET /strategies} serialises it as a JSON Schema; the React form
component renders a complete parameter UI from that schema with no per-strategy
frontend code. Adding the sixth-through-eleventh strategies (PR~\#7) cost zero
frontend work as a result \textemdash{} the single biggest productivity
multiplier in the project.

\paragraph{The six-agent abstraction.} Splitting the runtime into named agents,
each inheriting \texttt{BaseAgent[TIn, TOut]} with a single public
\texttt{run(payload)} method, gave us seam-by-seam testability and a natural
docking point for the LLM orchestrator. The boundary between the two LLM-backed
agents and the four deterministic agents was the right boundary: it kept the test
suite reproducible while leaving the LLM path open as an opt-in feature.

\paragraph{The $\{-1, 0, +1\}$ signal contract.} Forbidding strategies to manage
their own position size, commissions, or fills made the engine the single point of
truth for the bar-$t \to$ bar-$t{+}1$ rule, and meant that the
look-ahead-bias oracle test in \texttt{tests/test\_engine\_no\_lookahead.py} is the
only invariant guard we need. Every strategy passes through the same fill code
path; a bug anywhere in fill semantics surfaces everywhere at once.

\paragraph{Native trading calendars per asset class.} The temptation to reindex
every asset onto NYSE for ``consistency'' would have silently misstated the
annualisation factor for FX and crypto by 5\textendash{}50\%. Building
calendar-aware reindexing into the cleaner from the start (PR~\#2) cost a few
hours and saved us from a class of subtle correctness bugs.

\subsection{What did not work}

\paragraph{Yahoo Finance's 730-day hourly cap.} The cap is a hard external
constraint that propagates into the UI (date pickers had to gain a per-timeframe
maximum-look-back constraint) and into the data layer (the Crypto Fetcher learned
to fall back to Binance via \texttt{ccxt}). We absorbed the cap; the alternatives
(paid intraday data, scraping) cost more than the limitation imposes.

\paragraph{LLM tool-call JSON parsing.} The Orchestrator Agent parses tool calls
out of the model's response by extracting a JSON envelope. This is fragile:
models occasionally wrap the JSON in markdown fences or prose. We chose
provider-portability over robustness for v1; a future iteration could fall back
to provider-native function calling where available.

\paragraph{Asynchronous backtest execution.} The first design had submission
return immediately and the UI poll for status; we backed it out to a synchronous
\texttt{POST /backtests} because runs on daily bars complete in under a second
and the async machinery added an entire state-management layer that nothing in
v1 actually needed (Section~\ref{sec:diary}).

\subsection{What we would do differently}

\begin{itemize}
  \item \textbf{Walk-forward UI from day one.} The Strategy Agent already exposes a
        chronological train/test split, but the UI does not surface it. Without a UI
        affordance, users will tune parameters by re-running the full backtest
        \textemdash{} the canonical recipe for overfitting~\cite{bailey2014}. The
        UI for walk-forward should have been a v1 feature, not a v1.1 candidate.
  \item \textbf{Postgres as the default development database.} SQLite was fast to set
        up but does not catch the subset of bugs that only show up under concurrent
        writes or under strict referential integrity. The schema is already
        Postgres-compatible (SQLAlchemy generics only); next time we would default
        to Postgres and let users opt down to SQLite.
  \item \textbf{Earlier CI hygiene.} The \texttt{check-no-build-artifacts.yml}
        workflow should have been the first \texttt{.github/workflow} we wrote, not
        the third. Build-artefact leakage produced noisy diffs for a week before we
        formalised the guard.
  \item \textbf{Smaller, more frequent PRs.} The 32 development-cycle commits landed
        via only 8 merged PRs; some touched a dozen files across two subsystems.
        Reviewers (human or AI) work better on small, focused diffs.
\end{itemize}

\subsection{What carries forward}

The patterns we would copy verbatim into the next project: the
schema-as-API discipline, the agent abstraction with a single public
\texttt{run} method, the deterministic-by-default LLM stance with
\texttt{NullProvider} for tests, the per-asset-class trading-calendar layer, the
$\{-1, 0, +1\}$ signal contract, and the ``\emph{Last verified against code}''
footer convention on every documentation file.
